% !PRINTED 417
\weberpage{417}
Auch wie die allgemeine Gleichung $5^{\mathrm{ten}}$ Grades auf eine diesen Voraussetzungen entsprechende transformirt werden kann, ist dort gezeigt.

Unter der Voraussetzung (4) hat das Bestehen der Relationen (1) keinen Einfluss auf die Schlüsse des \S.\,108, und wir haben dann die folgenden Sätze.

Es sei $x_0, x_1, \ldots, x_{n-1}$ ein System von Variablen, das den Relationen $A=0$, $B=0$, \ldots\ unterworfen, aber sonst unabhängig ist, und es sei $\mathfrak{A}$ eine Permutationsgruppe der $x$, deren Permutationen die Functionen $A, B, \ldots$ ungeändert lassen.

\gesperrt{Soll die Gleichung, deren Wurzeln die Grössen $x$ sind, eine Resolvente}
\weq{6}{F(u,z) = 0}
\gesperrt{besitzen, die ausser numerischen Coëfficienten nur einen Parameter enthält, der eine durch die Gruppe $\mathfrak{A}$ ungeänderte Function der $x$ ist, so muss eine Function $\vartheta$ der Variablen $x_0, \ldots, x_{n-1}$ existiren, die durch die Permutationen $\alpha$ von $\mathfrak{A}$ (mit Rücksicht auf die Relationen $A=0$, $B=0$, \ldots) lineare Substitutionen erleidet}
\weq{7}{\vartheta_\alpha = \dfrac{a\vartheta+b}{c\vartheta+d},}
\gesperrt{worin $a, b, c, d$ numerische Coëfficienten sind. Die Wurzeln der Gleichung (6) sind rational durch $\vartheta$ ausdrückbar, und folglich ist auch $\vartheta$ Wurzel einer Gleichung mit einem Parameter}
\weq{8}{\Phi(\vartheta, z) = 0.}

Ist (6) eine Totalresolvente, was immer eintritt, wenn die Gruppe $\mathfrak{A}$ einfach ist, so können auch die $x$ rational durch $\vartheta$ und die zur Gruppe $\mathfrak{A}$ gehörigen Functionen ausgedrückt werden.

Bis hierher findet also vollständige Uebereinstimmung mit den Resultaten des \S.\,108 statt.

Die Beschränkung aber, auf die wir im \S.\,109 gestossen sind, wird hier nicht mehr nothwendig eintreten, weil jetzt nicht mehr, wie im Falle völlig unabhängiger Variablen, aus der Gleichheit zweier gebrochener Functionen auf die Uebereinstimmung von Zähler und Nenner geschlossen werden kann.

Die linearen Substitutionen (7) bilden eine mit $\mathfrak{A}$ (ein- oder mehrstufig) isomorphe Gruppe $P$, die mit einer unserer
%CONT>>
